% Show the retained closed intervals in the same style as interval halving.
% Develop the construction before naming its limiting set in the body.
\begin{tikzpicture}[interval diagram,x=6cm]
  \node[anchor=south west] at (-0.16,0.85) {Remaining intervals};

  \draw (0,0) -- (1,0);
  \fill (0,0) circle[radius=1.6pt];
  \fill (1,0) circle[radius=1.6pt];
  \node[left=6pt] at (0,0) {$E_0$};
  \node[below=5pt] at (0,0) {$0$};
  \node[below=5pt] at (1,0) {$1$};

  \foreach \a/\b in {0/1,2/3} {
    \draw ({\a/3},-1.8) -- ({\b/3},-1.8);
    \fill ({\a/3},-1.8) circle[radius=1.6pt];
    \fill ({\b/3},-1.8) circle[radius=1.6pt];
  }
  \node[left=6pt] at (0,-1.8) {$E_1$};
  \node[above=5pt] at (1/6,-1.8) {digit $0$};
  \node[above=5pt] at (5/6,-1.8) {digit $2$};
  \node[below=5pt] at (1/3,-1.8) {$1/3$};
  \node[below=5pt] at (2/3,-1.8) {$2/3$};

  \foreach \a in {0,2,6,8} {
    \draw ({\a/9},-3.8) -- ({(\a+1)/9},-3.8);
    \fill ({\a/9},-3.8) circle[radius=1.6pt];
    \fill ({(\a+1)/9},-3.8) circle[radius=1.6pt];
  }
  \node[left=6pt] at (0,-3.8) {$E_2$};

  \foreach \a in {0,2,6,8,18,20,24,26} {
    \draw ({\a/27},-5.3) -- ({(\a+1)/27},-5.3);
    \fill ({\a/27},-5.3) circle[radius=1.6pt];
    \fill ({(\a+1)/27},-5.3) circle[radius=1.6pt];
  }
  \node[left=6pt] at (0,-5.3) {$E_3$};
  \node at (0.5,-6.45) {$\vdots$};
\end{tikzpicture}
